\chapter{Recomendaciones}\label{cap:recomendaciones}
A partir de las limitaciones identificadas se formulan recomendaciones técnicas y académicas.

\section{Recomendaciones técnicas}
Estas extensiones conservan el principio de implementación desde cero.

\begin{enumerate}[leftmargin=2em,itemsep=1pt]
\item Añadir pruebas automatizadas al repositorio oficial, empezando por las operaciones de \texttt{Vector}, \texttt{MSE} y \texttt{Gradient}, y adaptar del motor complementario las pruebas de \emph{gradient checking}.
\item Corregir el error \texttt{clippy::modulo\_one} y eliminar o utilizar el código que hoy genera advertencias, como \texttt{Matrix}, \texttt{Network} y las funciones \texttt{random}.
\item Incorporar al repositorio oficial la retropropagación multicapa del motor complementario y un experimento XOR que justifique las capas ocultas.
\item Ampliar el experimento de vocales con variantes ruidosas y un conjunto de prueba separado, a matrices de mayor resolución o a todo el alfabeto.
\item Exportar los resultados de cada experimento a CSV directamente desde el programa y ofrecer una opción para desactivar las pausas de la demostración.
\item Explorar mini-lotes, optimizadores adaptativos como Momentum o Adam y regularización $L_2$ para conjuntos de datos más grandes.
\end{enumerate}

\section{Recomendaciones académicas}
El repositorio puede aprovecharse como recurso didáctico después de la feria.

\begin{enumerate}[leftmargin=2em,itemsep=1pt]
\item Usar el repositorio como laboratorio de Programación 1: reproducir con calculadora las actualizaciones impresas y modificar los parámetros.
\item Diseñar ejercicios graduados en los que los estudiantes implementen una nueva función de activación o de pérdida y la validen con \emph{gradient checking}.
\item Articular el proyecto con las asignaturas de Álgebra Lineal y Cálculo, presentando la retropropagación como aplicación directa de la regla de la cadena y del producto matricial.
\item Mantener el repositorio como proyecto abierto del Semillero de Investigación del IICCT, de modo que futuras cohortes puedan ampliarlo y documentar sus aportes.
\end{enumerate}
